\documentclass[10pt, letterpaper]{article}

\usepackage{../../../../template/template}

\setDocTitle{Verbale interno 31/03/2026}
\setDocVersion{-}
\setDocData{31/03/2026}
\setDocIndex{vi\_2026\_03\_31}
\setDocState{1}
\setDocRecipient{1}

\begin{document}

\makeTitlePage

% \newpage
% \addChangelog{0.1.0}{11/03/2026}{F. Pasqual}{C. Libralato}{\noOne}{\firstDraft}
% \makeChangelog

\newpage

\tableofcontents

\newpage

% -------------------------------------------------
% Informazioni riunione
% -------------------------------------------------
\makeMeetingInfoTable{31/03/2026}{14:30}{15:00}{via \textit{Discord$_G$}}

\addParticipant{Valeria}{Baleanu}{Programmatore}{P}
\addParticipant{Leonardo}{Pellizzon}{Programmatore}{P}
\addParticipant{Luca}{Granziero}{Amministratore}{P}
\addParticipant{Francesco}{Pasqual}{Programmatore}{P}
\addParticipant{Giuseppe}{De Fina}{Verificatore}{P}
\addParticipant{Christian}{Libralato}{Responsabile}{P}
\addParticipant{Filippo}{Venzo}{Programmatore}{P}
\makeMeetingParticipantsTable

\addPrecedentTodo{vi\_2026\_03\_24.a1}{\noOne}{\teamName}{31/03/2026}
\addPrecedentTodo{vi\_2026\_03\_24.a1}{\#24}{F. Venzo}{31/03/2026}
\addPrecedentTodo{vi\_2026\_03\_24.a1}{\noOne}{\teamName}{26/03/2026}
\addPrecedentTodo{vi\_2026\_03\_24.a1}{\#224}{C. Libralato}{31/03/2026}
\addPrecedentTodo{vi\_2026\_03\_24.a1}{\#225}{C. Libralato}{31/03/2026}

\makePrecedentTodoTable
% -------------------------------------------------
% Ordine del giorno
% -------------------------------------------------
% -------------------------------------------------
% Ordine del giorno
% -------------------------------------------------
\section{Ordine del giorno}
\begin{itemize}
    \item Stato avanzamento \textit{backend$_G$};
    \item aggiornamento e completamento della documentazione PDP e PDQ;
    \item coordinamento del lavoro sul branch di integrazione e gestione merge;
    \item dubbi.
  \end{itemize}

% -------------------------------------------------
% Approfondimento
% -------------------------------------------------
\section{Approfondimento}
\subsection*{Stato avanzamento \textit{backend}}
È stato discusso il completamento della parte di subscriptions sul \textit{backend}, incluso l'aggiornamento della cache 
e delle time series su variazione dei \textit{data points}. Sono state dichiarate terminate le modifiche delle analytics e l'inclusione dei suggerimenti lato \textit{backend}.

Il team ha sottolineato la necessità di salvare le notifiche generate dagli allarmi su una tabella nel database per
permettere la consultazione da parte dell'utente finale, considerando che il \textit{frontend$_G$} ha già una funzionalità per la visualizzazione delle notifiche.
La struttura della tabella è già stata modellata e necessita solo di essere implementata.

\subsection*{Aggiornamento e completamento della documentazione PDP e PDQ}
È emersa l'esigenza di aggiornare quanto prima la documentazione PDP e PDQ per lo sprint 10, con un controllo sui dati
e formule, in modo da poter poi procedere all'aggiornamento nello sprint 11. È stato sottolineato che gli aggiornamenti 
devono essere comunicati chiaramente al gruppo per monitorare i progressi.

\subsection*{Coordinamento del lavoro sul branch di integrazione e gestione merge}
È stato proposto di creare un branch chiamato 'integrazione' da mantenere stabile e con i test funzionanti, su cui tutti devono basare il proprio lavoro e dove devono essere fatte piccole modifiche testate.
È stata poi discussa la gestione dei branch aperti e la necessità di chiudere quelli non più usati, con l'obiettivo di migliorare la continuità e stabilità del codice.

\subsection*{Dubbi}
Sono emersi alcuni dubbi riguardo le seguenti tematiche:
\begin{itemize}
    \item Il grado di conformità al codice che i diagrammi UML devono possedere;
    \item il ruolo a cui è affidata la responsabilità di stesura dei test;
    \item il rispetto realistico delle scadenze.
\end{itemize} 
Le questioni saranno risolte tramite interazione con la Proponente e durante il diario di bordo.

% -------------------------------------------------
% Decisioni
% -------------------------------------------------
\addDecision{Salvataggio notifiche generate dagli allarmi}

\makeDecisionTable

% -------------------------------------------------
% Azioni
% -------------------------------------------------
\addTodo{\noOne}{Stesura verbale}{C. Libralato}{05/04/2026}
\addTodo{\noOne}{Aggiornamento PDP e PDQ}{L. Granziero}{05/04/2026}
\addTodo{\noOne}{Implementare salvataggio notifiche degli allarmi}{F. Venzo}{05/04/2026}
\addTodo{\noOne}{Chiudere i branch datati}{L. Granziero}{05/04/2026}

\makeTodoTable


\end{document}